\documentclass[11pt,a4paper]{article}

% -----------------------------------------------------------------------------
% Bamberger Modell III -- arXiv-style LaTeX main file
% -----------------------------------------------------------------------------

\usepackage[T1]{fontenc}
\usepackage[utf8]{inputenc}
\usepackage[ngerman]{babel}
\usepackage{lmodern}
\usepackage{microtype}

\usepackage{amsmath,amssymb,amsfonts,amsthm,mathtools}
\usepackage{geometry}
\usepackage{enumitem}
\usepackage{booktabs}
\usepackage{array}
\usepackage{csquotes}
\usepackage{hyperref}
\usepackage[nameinlink,capitalize,noabbrev]{cleveref}
\usepackage{graphicx}

\geometry{margin=2.7cm}

% Robust graphics setup
\DeclareGraphicsExtensions{.pdf,.png,.jpg,.jpeg}
\graphicspath{{./}{fig/}{figs/}{images/}{img/}}
\newcommand{\includegraphicssafe}[2][]{%
  \IfFileExists{#2}{\includegraphics[#1]{#2}}{%
    \fbox{\parbox{0.85\linewidth}{\centering Fehlende Grafik: \texttt{#2}}}%
  }%
}

% Theorem environments -- deutsch benannt
\newtheorem{satz}{Satz}[section]
\newtheorem{lemma}[satz]{Lemma}
\newtheorem{korollar}[satz]{Korollar}
\theoremstyle{definition}
\newtheorem{definition}[satz]{Definition}
\newtheorem{hypothese}[satz]{Hypothese}
\theoremstyle{remark}
\newtheorem{bemerkung}[satz]{Bemerkung}
\newtheorem{beispiel}[satz]{Beispiel}

% Operators and shortcuts
\DeclareMathOperator{\rad}{rad}
\DeclareMathOperator{\sgn}{sgn}
\DeclareMathOperator{\Area}{Area}
\DeclareMathOperator{\Span}{span}
\newcommand{\N}{\mathbb N}
\newcommand{\Z}{\mathbb Z}
\newcommand{\R}{\mathbb R}
\newcommand{\C}{\mathbb C}
\newcommand{\HH}{\mathbb H}
\newcommand{\eps}{\varepsilon}
\newcommand{\phig}{\varphi}
\newcommand{\eabc}{EABC}
\newcommand{\vect}[1]{\mathbf{#1}}

\title{Bamberger Modell III\\
Tetraedrische Spin-Geometrie, quaternionische Bindungen,\\
ikosaedrische Bahnimpulse und modulare Phasenstrukturen\\
auf den natuerlichen Zahlen}

\author{Thomas Hoffbauer\\Bamberg, Deutschland}
\date{\today}

\begin{document}

\maketitle

\begin{abstract}
Wir entwickeln eine arithmetisch-geometrische Struktur auf den natuerlichen Zahlen,
die aus den invertierbaren Restklassen modulo 12, der Quaternionenalgebra, den
regulaeren Polyedern, Morley-Dreiecken und modularen Formen hervorgeht. Im Zentrum
steht die Beobachtung, dass die vier Einheiten
\[
(\Z/12\Z)^\times=\{1,5,7,11\}
\]
eine natuerliche tetraedrische Geometrie erzeugen. Diese wird mit Hamiltons
Quaternionen, den Hurwitz-Einheiten, Dedekinds Eta-Funktion, den
Rogers--Ramanujan-Identitaeten und der Ikosaedergeometrie Felix Kleins verknuepft.

Das resultierende Modell trennt eine Zahl in Skala, Spin, Bindung, Orbit und
Spinorphase. Daraus entsteht eine arithmetische Teilchenspektroskopie aus
elementaren Spinoren, Mesonen, Baryonen und Singuletts. Die physikalischen
Begriffe werden dabei als Modellanalogien eingefuehrt; die mathematische Basis
besteht aus Restklassen, Gruppen, Tetraedergeometrie, Quaternionen und modularen
Phasen.
\end{abstract}

\tableofcontents

\section{Einleitung}

Die klassische Zahlentheorie beschreibt natuerliche Zahlen vor allem durch ihre
Groesse, ihre Teilbarkeit und ihre Primfaktorzerlegung. Das vorliegende Modell
untersucht eine zusaetzliche Hypothese: Die Primfaktorisierung traegt eine
geometrische Orientierungsstruktur. Diese Struktur beginnt bei den vier
invertierbaren Restklassen modulo 12,
\[
1,5,7,11,
\]
und fuehrt ueber eine tetraedrische Geometrie zu Spin-, Bindungs- und
Bahnimpulsanalogien.

Der Begriff \enquote{Bamberger Modell III} bezeichnet in diesem Text keine
abgeschlossene physikalische Theorie, sondern ein mathematisch strukturiertes
Forschungsprogramm. Seine sichere Ebene ist die Algebra und Geometrie der
Restklassen. Seine spekulative Ebene ist die Interpretation dieser Geometrie als
arithmetische Vorform von Spin, Kern, Huelle, Bindung und Masse.

\section{Historische Wurzeln}

\subsection{Euler und die Primzahlen}

Die moderne analytische Zahlentheorie beginnt wesentlich mit Eulers
Produktdarstellung der Zetafunktion:
\[
\zeta(s)=\prod_p \frac{1}{1-p^{-s}}.
\]
Sie macht deutlich, dass Primzahlen die multiplikativen Atome der natuerlichen
Zahlen sind. Im Bamberger Modell bleiben Primzahlen atomare Traeger, erhalten
aber zusaetzlich eine Orientierung durch ihre Restklasse modulo 12.

\subsection{Hamilton und die Quaternionen}

Hamiltons Quaternionenalgebra
\[
\HH=\Span_{\R}\{1,i,j,k\}
\]
mit
\[
i^2=j^2=k^2=ijk=-1
\]
ist die klassische nichtkommutative Erweiterung der komplexen Zahlen. Sie ist
fundamental fuer Rotationen und Spinoren, da die normierten Quaternionen eine
Realisierung von \(SU(2)\) liefern. Diese Tatsache bildet den historischen und
mathematischen Hintergrund fuer die spaetere Hebung der kommutativen
\(EABC\)-Algebra in eine orientierte Quaternionenstruktur.

\subsection{Hurwitz und die 24 Einheiten}

Der Ring der Hurwitz-Quaternionen besitzt 24 Einheiten. Diese bilden die
Scheitelpunkte des regulaeren 24-Zellers. Im Modell werden diese 24 Einheiten
als geometrische Entsprechung einer modularen 24-Phase gelesen. Sie verbinden
die Quaternionengeometrie mit der Eta-Struktur.

\subsection{Dedekind und die Eta-Funktion}

Die Dedekindsche Eta-Funktion lautet
\[
\eta(\tau)=q^{1/24}\prod_{n\ge 1}(1-q^n),\qquad q=e^{2\pi i\tau}.
\]
Sie besitzt die Transformation
\[
\eta(\tau+1)=e^{\pi i/12}\eta(\tau).
\]
Damit treten die Zahlen 12 und 24 auf natuerliche Weise auf. Im Modell liefert
\(\eta\) die modulare Phasenhuelle ueber der sichtbaren Restklassenstruktur.

\subsection{Ramanujan und die modulare Resonanz}

Ramanujan fand tiefe Identitaeten zwischen modularen Funktionen,
Kettenbruechen und speziellen Exponentialwerten. Besonders relevant ist der
Rogers--Ramanujan-Kettenbruch
\[
R(q)=\frac{q^{1/5}}{1+\frac{q}{1+\frac{q^2}{1+\cdots}}}.
\]
Fuer Werte wie \(q=e^{-\pi}\) und \(q=e^{-2\pi}\) treten geschlossene Formen auf,
die mit dem Goldenen Schnitt
\[
\phig=\frac{1+\sqrt5}{2}
\]
verknuepft sind. Im Modell wird diese Beziehung als Bruecke zwischen modularer
Skala und ikosaedrischer Geometrie interpretiert.

\subsection{Klein und das Ikosaeder}

Felix Klein erkannte die zentrale Rolle des Ikosaeders und seiner
Rotationsgruppe \(A_5\) fuer die Theorie modularer Funktionen. Da der Goldene
Schnitt in den Koordinaten des Ikosaeders unvermeidlich erscheint, wird der
Uebergang vom Tetraeder zum Ikosaeder im Bamberger Modell als goldene
Verfeinerung der urspruenglichen Spin-Geometrie aufgefasst.

\section{Die EABC-Algebra}

\begin{definition}[EABC-Klassen]
Die vier invertierbaren Restklassen modulo 12 werden bezeichnet durch
\[
E\equiv 1,\qquad A\equiv 5,\qquad B\equiv 7,\qquad C\equiv 11\pmod{12}.
\]
\end{definition}

Die Multiplikation modulo 12 ergibt
\[
A^2=B^2=C^2=E,
\]
und
\[
AB=C,\qquad BC=A,\qquad CA=B.
\]
Diese Struktur ist isomorph zur Klein-Vierergruppe \(V_4\).

\begin{bemerkung}
Die \(EABC\)-Algebra ist zunaechst kommutativ. Sie beschreibt daher noch keine
orientierte oder fermionische Struktur. Eine solche entsteht erst durch die
quaternionische Hebung.
\end{bemerkung}

\section{Der tetraedrische Spinraum}

Wir ordnen den vier \(EABC\)-Klassen die Vektoren
\[
E=(1,1,1),
\]
\[
A=(-1,-1,1),
\]
\[
B=(1,-1,-1),
\]
\[
C=(-1,1,-1)
\]
zu.

\begin{satz}[Tetraedersatz]
Die vier Vektoren \(E,A,B,C\) bilden die Ecken eines regulaeren Tetraeders.
\end{satz}

\begin{proof}
Fuer jedes verschiedene Paar \(X,Y\in\{E,A,B,C\}\) gilt
\[
\|X-Y\|=2\sqrt2.
\]
Alle sechs Kanten haben also dieselbe Laenge. Damit bilden die vier Punkte ein
regulaeres Tetraeder.
\end{proof}

Zudem gilt
\[
E+A+B+C=0.
\]
Der vollstaendige Vierer ist daher spinneutral.

\begin{definition}[EABC-Spin]
Der tetraedrische Spin eines Halbklein-Zustands \(H\in\{E,A,B,C\}\) ist
\[
S(H)=H\in\R^3.
\]
\end{definition}

\section{Primzahlen als elementare Spinoren}

Jede Primzahl \(p>3\) liegt in genau einer der vier Klassen \(E,A,B,C\). Damit
traegt jede solche Primzahl einen elementaren Tetraederspin. Diese Sichtweise
motiviert die Bezeichnung \emph{arithmetischer Spinor} fuer eine einzelne
quaternionisch reine Primklasse.

\begin{definition}[Radix-Besetzung]
Fuer eine natuerliche Zahl \(n\) sei
\[
\rad(n)=\prod_{p\mid n}p.
\]
Wir schreiben \((r_E,r_A,r_B,r_C)\) fuer die Anzahl verschiedener Primfaktoren
\(p>3\) von \(n\) in den Klassen \(E,A,B,C\).
\end{definition}

Der Radix-Spin ist
\[
S_{\rad}(n)=r_EE+r_AA+r_BB+r_CC.
\]

\section{Mesonen, Baryonen und Singuletts}

Die niedrigsten Kombinationsklassen ergeben eine arithmetische
Teilchenspektroskopie.

\subsection{Einer: Spinoren}

Ein einzelner Zustand \(X\in\{E,A,B,C\}\) besitzt
\[
\|X\|^2=3.
\]
Er ist nicht neutral und wird als elementarer Spinor interpretiert.

\subsection{Zweier: Mesonen}

Zweierkombinationen erzeugen Achsenzustande:
\[
E+A=(0,0,2),\qquad E+B=(2,0,0),\qquad E+C=(0,2,0),
\]
\[
A+B=(0,-2,0),\qquad A+C=(-2,0,0),\qquad B+C=(0,0,-2).
\]
Diese Paare werden als arithmetische Mesonen oder Dipole bezeichnet.

\subsection{Dreier: Baryonen}

Dreierkombinationen sind Lochzustaende:
\[
A+B+C=-E,
\]
\[
E+B+C=-A,
\]
\[
E+A+C=-B,
\]
\[
E+A+B=-C.
\]
Sie tragen die Gegenrichtung des fehlenden vierten Zustands.

\subsection{Vierer: Singuletts}

Der vollstaendige Vierer erfuellt
\[
E+A+B+C=0.
\]
Er ist das arithmetische Singulett oder die geschlossene Schale des Modells.

\section{Kern-Huelle-Struktur}

Die numerischen und algebraischen Befunde legen folgende Kern-Huelle-Lesart nahe:
\[
E=\text{Kernkanal},
\]
\[
B,C=\text{Huellendpolaritaeten},
\]
\[
A=\text{Bindungskanal}.
\]
Der elementare Huellendipol ist
\[
D_{BC}=B+C=(0,0,-2).
\]
Seine Gegenorientierung ist
\[
-D_{BC}=(0,0,2).
\]

\begin{hypothese}[Kern-Huelle-Regel]
Der Zustand \(E\) beschreibt den Kernkanal, das Paar \(B+C\) den elementaren
Huellendipol und \(A\) den Bindungskanal zwischen Kern und Huelle.
\end{hypothese}

\section{Quaternionische Hebung und fermionische Vorzeichen}

Die sichtbare \(EABC\)-Algebra ist kommutativ. Zur Orientierung wird sie gehoben
nach
\[
E\mapsto 1,
\qquad
A\mapsto i,
\qquad
B\mapsto j,
\qquad
C\mapsto k.
\]
Dann gilt
\[
ij=k,
\qquad
jk=i,
\qquad
ki=j,
\]
aber
\[
ji=-k,
\qquad
kj=-i,
\qquad
ik=-j.
\]
Insbesondere:
\[
BC=+A,
\qquad
CB=-A.
\]

\begin{satz}[Fermionische Hebung]
Die kommutative \(V_4\)-Struktur der sichtbaren \(EABC\)-Klassen wird durch die
quaternionische Hebung zu einer orientierten \(Q_8\)-Struktur erweitert. Die
Vertauschung zweier Huellenkomponenten erzeugt einen Vorzeichenwechsel.
\end{satz}

\section{Eta-24-Phase}

Die Dedekindsche Eta-Funktion erzeugt durch
\[
\eta(\tau+1)=e^{\pi i/12}\eta(\tau)
\]
eine natuerliche 24-periodische Phase. Wir setzen daher
\[
\eps\in\Z/24\Z.
\]
Fuer den orientierten Bindungskanal wird definiert:
\[
+A\leftrightarrow \eps=0,
\qquad
-A\leftrightarrow \eps=12\pmod{24}.
\]
Die Vertauschung
\[
BC\longrightarrow CB
\]
entspricht dem Halbumlauf
\[
\eps\longrightarrow \eps+12.
\]
Nach zwei Vertauschungen gilt
\[
12+12=24\equiv 0.
\]
Damit entsteht die spinorische Signatur
\[
360^\circ\mapsto -1,
\qquad
720^\circ\mapsto +1.
\]

\section{Morley-Kopplung}

Fuer jede Dreiecksflache des Spin-Tetraeders kann das Morley-Dreieck gebildet
werden. Der Morley-Skalar sei
\[
\mu_M=\frac{\Area(M)}{\Area(F)},
\]
wobei \(F\) die Ausgangsflaeche und \(M\) das Morley-Dreieck bezeichnet. In den
bisherigen numerischen Modellrechnungen wird
\[
\mu_M\approx 0.03414827948
\]
verwendet.

Der Morley-Skalar koppelt Spin- und Orbitstruktur:
\[
J=S+\mu_M L.
\]

\section{Goldene Verfeinerung und Ikosaeder}

Der Goldene Schnitt
\[
\phig=\frac{1+\sqrt5}{2}
\]
erscheint nicht im Tetraeder selbst, sondern beim Uebergang zur
Ikosaedergeometrie. Die Ecken des Ikosaeders koennen dargestellt werden als
\[
(0,\pm1,\pm\phig),
\]
\[
(\pm1,\pm\phig,0),
\]
\[
(\pm\phig,0,\pm1).
\]
Die Rotationsgruppe des Ikosaeders ist \(A_5\). Im Modell wird der Uebergang
\[
A_4\longrightarrow A_5
\]
als goldene Verfeinerung der tetraedrischen Spinstruktur interpretiert.

\section{Bahnimpuls aus geschlossenen Umlaeufen}

Fuer einen geschlossenen Umlauf
\[
P_1\to P_2\to\cdots\to P_m\to P_1
\]
wird der diskrete Bahnimpuls definiert durch
\[
L=\sum_{r=1}^{m}P_r\times P_{r+1}.
\]
Der Spin sitzt auf den Punkten des Tetraeders; der Bahnimpuls entsteht erst
durch gerichtete Umlaeufe.

\begin{bemerkung}
Auf dem Tetraeder erzeugen die Dreiecksumlaeufe das duale Tetraeder. Die
goldene Ikosaeder-Verfeinerung fuehrt zusaetzlich zu Fuenferumlaeufen und damit
zu einer natuerlichen \(A_5\)-Orbitstruktur.
\end{bemerkung}

\section{Masse als arithmetische Rigiditaet}

Masse wird im Modell nicht als Substanzmenge definiert, sondern als Rigiditaet
einer gekoppelten Spin-Orbit-Bindungsstruktur.

Eine erste Modellformel lautet
\[
m(n)=\Lambda(G,\eta)\left[\alpha\Omega(V)+\beta\|S_{\rad}(n)\|^2+
\gamma\mu_M\|L_{\phig}(n)\|^2+\delta B_{\rm ary}(n)\right].
\]
Hier bezeichnet \(\Omega(V)\) die Faktorisierungstiefe, \(S_{\rad}\) den
Radix-Spin, \(L_{\phig}\) den goldverfeinerten Bahnimpuls und \(B_{\rm ary}\)
einen baryonischen Dreierbeitrag.

\begin{hypothese}[Masse als Rigiditaet]
Masse misst die eingefrorene Stabilitaet einer arithmetischen Struktur. Sie ist
nicht identisch mit Spin, sondern entsteht aus Spinrest, Orbit, Bindung,
Faktorisierungstiefe und modularer Phase.
\end{hypothese}

\section{Die vollstaendige Zerlegung einer Zahl}

Eine natuerliche Zahl wird modellhaft beschrieben durch
\[
n=(G,H,V,\eps).
\]
Dabei bezeichnet
\[
G=2^a3^b
\]
den glatten Skalenanteil,
\[
H\in\{E,A,B,C\}
\]
den halbkleinen Spinanteil,
\[
V
\]
den vollkleinen Bindungs- und Orbitanteil, und
\[
\eps\in\Z/24\Z
\]
die modulare Spinorphase.

\section{Hauptsatz des Bamberger Modells III}

\begin{satz}[Struktur der natuerlichen Zahlen]
Im Bamberger Modell III besitzt jede natuerliche Zahl eine vierstufige Struktur
\[
n=(G,H,V,\eps).
\]
Der glatte Anteil \(G\) liefert die Skala. Der halbkleine Anteil \(H\) liefert
den tetraedrischen Spin. Der vollkleine Anteil \(V\) liefert Bindung und Orbit.
Die Eta-Phase \(\eps\) speichert die spinorische Orientierung.
\end{satz}

\begin{korollar}
Die sichtbare Restklassenstruktur wird durch \(V_4\) beschrieben, die
fermionische Orientierung durch die quaternionische Hebung \(Q_8\), und die
modulare Spinorphase durch \(\Z/24\Z\).
\end{korollar}

\section{Offene Probleme}

Die folgenden Punkte bleiben offen:
\begin{enumerate}[label=\arabic*.]
\item eine exakte Herleitung physikalischer Massen aus der Rigiditaetsformel;
\item die volle fermionische Antikommutationsalgebra;
\item die Praezisierung des Zusammenhangs zu \(SU(2)\) und \(SU(3)\);
\item die Einbettung in Hecke-Operatoren und modulare Formen;
\item die Beziehung zu Riemann-Nullstellen und spektralen Trace-Formeln;
\item die Rolle von Hurwitz-, 24-Zeller-, E8- und Leech-Strukturen;
\item numerische Tests der Radix-Spektren bis grosse Schranken.
\end{enumerate}

\section{Schluss}

Das Bamberger Modell III formuliert eine mehrschichtige Geometrie der
natuerlichen Zahlen. Die Restklassen modulo 12 erzeugen ein Tetraeder. Die
Quaternionen erzeugen Orientierung. Die Eta-Funktion erzeugt eine 24-Phase.
Ramanujansche Identitaeten verbinden modulare Skalen mit dem Goldenen Schnitt.
Die Ikosaedergeometrie liefert eine goldene Orbitverfeinerung. Daraus entsteht
ein Modell von Skala, Spin, Bindung, Orbit, Masse und Phase auf der
Faktorisierungsstruktur der natuerlichen Zahlen.

\begin{thebibliography}{99}

\bibitem{Euler}
L. Euler,
\emph{Introductio in analysin infinitorum}, 1748.

\bibitem{Hamilton}
W. R. Hamilton,
\emph{On Quaternions}, Proceedings of the Royal Irish Academy, 1844.

\bibitem{Hurwitz}
A. Hurwitz,
\emph{Ueber die Zahlentheorie der Quaternionen}, Nachrichten von der Gesellschaft der Wissenschaften zu Goettingen, 1896.

\bibitem{Dedekind}
R. Dedekind,
\emph{Erlaeuterungen zu den Fragmenten XXVIII}, in: Bernhard Riemann's gesammelte mathematische Werke, 1876.

\bibitem{Ramanujan}
S. Ramanujan,
\emph{Modular Equations and Approximations to \(\pi\)}, Quarterly Journal of Mathematics, 1914.

\bibitem{Klein}
F. Klein,
\emph{Vorlesungen ueber das Ikosaeder und die Aufloesung der Gleichungen vom fuenften Grade}, Teubner, 1884.

\bibitem{ConwaySmith}
J. H. Conway and D. A. Smith,
\emph{On Quaternions and Octonions}, A K Peters, 2003.

\bibitem{Serre}
J.-P. Serre,
\emph{A Course in Arithmetic}, Springer, 1973.

\bibitem{Apostol}
T. M. Apostol,
\emph{Modular Functions and Dirichlet Series in Number Theory}, Springer, 1990.

\end{thebibliography}

\end{document}
